We define the \initializeCalleeContextName{} constraint system as follows, using the \initializeContextName{} constraints from section~(\ref{hub: consistencies: context: constraints}),
\[
	\resizebox{.95\linewidth}{!}{$
	\left\{ \begin{array}{l}
		\initializeCalleeContext {
			anchorRow                = i                                      ,
			relOffset                = \relof                                 ,
			initializationParameters = \childFrameInitializationParameters ,
		} \vspace {4mm} \\
		\qquad \define
		\initializeContext{
			anchorRow                   = i                                                                                             ,
			relOffset                   = \relof                                                                                        ,
			contextNumber               = \cn\new_{i}                                                                                   ,
			callStackDepth              = 1 + \locCurrentCallStackDepth                                                                 ,
			isRoot                      = 0                                                                                             ,
			isStatic                    = (\textcolor{draculaorange}{\bigstar}) \quad \locChildContextIsStatic                          ,
			accountAddressHigh          = (\textcolor{draculaorange}{\bigstar}) \quad \locChildContextAccountAddressHi                  ,
			accountAddressLow           = (\textcolor{draculaorange}{\bigstar}) \quad \locChildContextAccountAddressLo                  ,
			accountDeploymentNumber     = (\textcolor{draculaorange}{\bigstar}) \quad \locChildContextAccountDeploymentNumber           ,
			byteCodeAddressHi           = \locDelegateOrCalleeAddressHi                                                                 ,
			byteCodeAddressLo           = \locDelegateOrCalleeAddressLo                                                                 ,
			byteCodeDeploymentNumber    = \locDelegateOrCalleeDeploymentNumber                                                          ,
			byteCodeDeploymentStatus    = \locDelegateOrCalleeDeploymentStatus                                                          ,
			byteCodeCodeFragmentIndex   = \locDelegateOrCalleeCfi                                                                       ,
			callerAddressHi             = (\textcolor{draculaorange}{\bigstar}) \quad \locChildContextCallerAddressHi                   ,
			callerAddressLo             = (\textcolor{draculaorange}{\bigstar}) \quad \locChildContextCallerAddressLo                   ,
			callValue                   = (\textcolor{draculaorange}{\bigstar}) \quad \locChildContextCallValue                         ,
			callDataContextNumber       = \cn_{i}                                                                                       ,
			callDataOffset              = \locTypeSafeCdo                                                                               ,
			callDataSize                = \locTypeSafeCds                                                                               ,
			returnAtOffset              = \locTypeSafeRao                                                                               ,
			returnAtCapacity            = \locTypeSafeRac                                                                               ,
		}
	\end{array} \right.
	$}
\]
All shorthands above were previously defined with the exception of those marked with $(\textcolor{draculaorange}{\bigstar})$
\begin{multicols}{2}
	\begin{itemize}
		\item \locChildContextAccountAddressHi
		\item \locChildContextAccountAddressLo
		\item \locChildContextAccountDeploymentNumber
		\item[\vspace{\fill}]
		\item \locChildContextCallerAddressHi
		\item \locChildContextCallerAddressLo
		\item \locChildContextIsStatic
		\item \locChildContextCallValue
	\end{itemize}
\end{multicols}
\noindent These shorthands will be defined in section~(\ref{hub: instruction handling: call: specialized: callee context parameters definition}).
